% Chapter 5 — Logical Design and Normalisation

\chapter{Logical Design and Normalisation}
\label{ch:logical}

\section{From drawing to tables}
Conceptual design answers ``what exists?'' Logical design answers ``how do we split information into tables so updates stay safe and queries stay clear?'' That process is \textbf{normalisation}: remove repeated groups, move repeating lists to their own tables, and connect rows with foreign keys.

\section{Why we use small lookup tables}
Instead of storing the word ``skip'' in every row of \texttt{activity\_log}, we store a tiny integer \texttt{type\_id} that points to \texttt{activity\_types}. Benefits:

\begin{itemize}[leftmargin=*]
  \item If we ever rename a label, we update \textbf{one row} in the lookup table.
  \item Typos cannot create fake categories---only listed types exist.
  \item Storage is smaller when millions of events accumulate.
\end{itemize}

\section{Functional dependencies (informal)}
\begin{itemize}[leftmargin=*]
  \item One \texttt{user\_id} determines current profile fields (name, email, role), with \texttt{email} separately unique.
  \item One \texttt{course\_id} determines its instructor and metadata (title, category), when not deleted.
  \item The pair (\texttt{user\_id}, \texttt{course\_id}) should identify \textbf{at most one} enrollment---enforced by a UNIQUE constraint.
\end{itemize}

\section{Why quiz answers live in their own table}
We could imagine cramming all answers into the attempt row as JSON. Instead, \texttt{quiz\_answers} stores one row per question per attempt. That mirrors the relational model used in textbooks: repeating groups become child rows. It also makes future reporting easier (``which option was chosen most often?'').

\section{Referential actions (delete behaviour)}
\textbf{Cascade} on user-owned child rows (sessions, answers) means removing a user removes dependent junk cleanly---important for privacy requests.

\textbf{Restrict} on lookup tables ensures we do not delete ``active'' status while enrollments still reference it.

\section{Relational Data Model Diagram (RDM)}
\label{sec:rdm-fig}

The RDM shows each \textbf{table name}, \textbf{primary key}, \textbf{foreign keys}, and major columns. Materialized views (\texttt{mv\_*}) can appear dashed or in a separate ``analytics'' panel because they are \textbf{derived}, not base tables.

\textbf{Insert your file} as \texttt{figures/rdm.pdf} or \texttt{figures/rdm.png}; uncomment \texttt{\textbackslash includegraphics} below.

\begin{figure}[htbp]
  \centering
  % \includegraphics[width=0.98\textwidth,height=0.62\textheight,keepaspectratio]{rdm.pdf}
  % \includegraphics[width=0.98\textwidth,height=0.62\textheight,keepaspectratio]{rdm.png}
  \fcolorbox{black!50}{black!2}{%
    \parbox{0.94\textwidth}{%
      \centering\vspace{3.6cm}
      {\large\textbf{Placeholder: Relational schema (RDM)}}\\[1.2em]
      \textit{Add file} \texttt{latex/figures/rdm.pdf} \textit{or} \texttt{rdm.png}\\[0.6em]
      \small Tables with PK/FK arrows; include lookups and quiz subgraph
      \vspace{3.6cm}
    }%
  }
  \caption{Relational data model for LearnTrack (replace placeholder with your diagram).}
  \label{fig:rdm}
\end{figure}

\section{Normalisation takeaway for presentations}
Say this in one breath: ``We separated identities, courses, enrollments, events, progress, and quiz structure into different tables so each fact lives in one place, connected by keys.'' That satisfies most DBMS oral expectations without jargon overload.
